\documentclass[11pt]{article}
\usepackage[T1]{fontenc}
\usepackage[utf8]{inputenc}
\usepackage{lmodern}
\usepackage[margin=0.9in,headheight=14pt]{geometry}
\usepackage{amsmath,amssymb}
\usepackage{booktabs,tabularx}
\usepackage[dvipsnames]{xcolor}
\usepackage{enumitem}
\usepackage{fancyhdr}
\usepackage{hyperref}
\usepackage{microtype}

\definecolor{navy}{HTML}{17365D}
\definecolor{teal}{HTML}{137C8B}
\definecolor{softblue}{HTML}{EAF1F8}
\definecolor{softgray}{HTML}{F3F5F7}
\definecolor{amber}{HTML}{8A5A00}
\hypersetup{colorlinks=true,linkcolor=navy,urlcolor=teal,
  pdftitle={MP1994V2 Proof Ledger}}
\pagestyle{fancy}
\fancyhf{}
\fancyhead[L]{\small MP1994V2}
\fancyhead[R]{\small Cumulative proof ledger}
\fancyfoot[C]{\thepage}
\setlength{\parindent}{0pt}
\setlength{\parskip}{0.55em}
\setlist{nosep,leftmargin=1.5em}
\allowdisplaybreaks
\newcommand{\code}[1]{\texttt{\detokenize{#1}}}
\newcommand{\completeGreen}{\textcolor{green!50!black}{\textbf{COMPLETE -- GREEN}}}

\title{\vspace{-2em}\textcolor{navy}{\textbf{MP1994V2 Proof Ledger}}}
\author{Mortensen--Pissarides (1994) Lean replication}
\date{29 July 2026}

\begin{document}
\maketitle

\colorbox{softblue}{\parbox{0.94\linewidth}{
\textbf{Purpose.} The milestone tracker records project state. This cumulative
ledger records exactly what Lean has proved, the hypotheses used, and the
economic meaning of each theorem.
}}

\section{M1: Surplus Bellman equation}

\begin{tabularx}{\linewidth}{@{}>{\bfseries}p{0.25\linewidth}X@{}}
Paper reference & Section 3, journal page 401, equation (8)\\
Adequacy status & \completeGreen\\
Lean source & \nolinkurl{MP1994V2/SteadyState/Surplus.lean}\\
\end{tabularx}

\subsection{Mathematical statement}

For the actual equilibrium surplus
\[
S(\varepsilon)=J(\varepsilon)+W(\varepsilon)-U,
\]
Equation (8) is a conditional representation theorem: for every object
satisfying \code{ValueEquilibrium P}, Lean proves that the surplus derived from
its value functions satisfies, for every idiosyncratic state \(\varepsilon\),
\begin{equation}
(r+\lambda)S(\varepsilon)
=p+\sigma\varepsilon-b
+\lambda\int \max\{S(x),0\}\,dF(x)
-\beta\theta q(\theta)S(\varepsilon_u).
\tag{8}
\end{equation}
Here \code{P.shock} is the probability law of the idiosyncratic shock; its
induced CDF is \(F\), and integration with respect to \(dF\) is represented in
Lean by integration against \code{P.shock}. Also, \(\theta q(\theta)\) is
\code{P.workerMeetingRate E.theta}. M1 does not prove existence of a
\code{ValueEquilibrium}.

\subsection{Economic interpretation}

\begin{itemize}
\item \(p+\sigma\varepsilon\): current match output.
\item \(b\): unemployment opportunity cost.
\item \(\lambda\int\max\{S(x),0\}\,dF(x)\): option value of future
idiosyncratic draws, allowing rejection of nonpositive continuation surplus.
\item \(\beta\theta q(\theta)S(\varepsilon_u)\): the employed worker's
foregone search gain at the upper-support job.
\end{itemize}

\subsection{Exact Lean declarations}

\begin{tabularx}{\linewidth}{@{}p{0.29\linewidth}X@{}}
\toprule
\textbf{Role} & \textbf{Declaration}\\
\midrule
Continuation decomposition &
\code{MP1994V2.continuationTerm_eq_integral_activeSurplus_sub}\\
Firm surplus share & \code{MP1994V2.ValueEquilibrium.firm_share}\\
Algebraic surplus flow &
\code{MP1994V2.ValueEquilibrium.surplus_flow_equation}\\
Minimal probability theorem &
\code{MP1994V2.ValueEquilibrium.surplus_bellman_of_probability}\\
Paper-level wrapper &
\code{MP1994V2.ValueEquilibrium.surplus_bellman}\\
\bottomrule
\end{tabularx}

\subsection{Assumptions}

\textbf{Used:}
\begin{itemize}
\item equation (3), through the definition of actual surplus;
\item equation (4), only at the upper state;
\item equations (5) and (6), added to cancel the wage;
\item equation (7), subtracted to account for unemployment value;
\item probability normalization of \code{P.shock};
\item \code{ValueEquilibrium.active_surplus_integrable}.
\end{itemize}

The minimal theorem assumes only
\code{[IsProbabilityMeasure P.shock]}. The wrapper takes
\code{D : ShockAssumptions P}, installs \code{D.isProbability} locally, and
uses no other field of \(D\).

\textbf{Not used:} parameter signs; \code{CoreEconomicAssumptions};
upper-support restrictions or maximality; atomlessness; mean zero; unit
variance; \code{ShockNormalizationAssumptions}; matching positivity or
monotonicity; \code{MatchingAssumptions}; equations (1) and (2).

\subsection{Readable proof}

Write the shared continuation term as
\[
C(\varepsilon)
=\int\left[\max\{S(x),0\}-S(\varepsilon)\right]\,dF(x).
\]
Adding equations (5) and (6) cancels the wage and sums the Nash weights:
\[
r[J(\varepsilon)+W(\varepsilon)]
=p+\sigma\varepsilon+\lambda C(\varepsilon).
\]
Subtracting equation (7), and using equation (3), gives
\[
rS(\varepsilon)
=p+\sigma\varepsilon-b+\lambda C(\varepsilon)
-\theta q(\theta)[W(\varepsilon_u)-U].
\]
Equation (4) at \(\varepsilon_u\) gives
\[
W(\varepsilon_u)-U=\beta S(\varepsilon_u).
\]
Probability normalization of \code{P.shock} and active-surplus integrability give
\[
C(\varepsilon)
=\int\max\{S(x),0\}\,dF(x)-S(\varepsilon).
\]
Substitution and moving \(-\lambda S(\varepsilon)\) to the left proves (8).

\subsection{Adequacy and next dependency}

The actual \code{ValueEquilibrium} surplus is used. Equation (8) is derived,
not added as an assumption or field. No cutoff, reservation rule, or affine
surplus result is used. It is conditional on an existing
\code{ValueEquilibrium} and does not establish that such an equilibrium
exists.

\textbf{Human-review conclusion.} Equation (8) matches the paper and is
derived non-circularly from equations (3)--(7). The actual equilibrium surplus
is used, with no cutoff, affinity result, or reduced equilibrium. The theorem
is explicitly conditional on an existing \code{ValueEquilibrium}.

M2 may subtract (8) at two productivity states to derive the surplus-difference
identity, then prove affine surplus and monotonicity consequences. No M2 result
is implemented here.

\vfill
\textit{Review space:}\\[4em]
\hrule

\end{document}
